\section{Empiece aquí: ruta autosuficiente de estudio}
\label{sec:programa-estudio-bibliografico}

Este volumen físico contiene el itinerario obligatorio: definiciones,
hipótesis, demostraciones y deducciones necesarias, metodologías, contratos
numéricos, resultados, tablas, discusión, ejercicios y comprobaciones. Puede
estudiarse de principio a fin sin abrir otro PDF. Siga las etapas en el orden
indicado y use las referencias bibliográficas únicamente para ampliar una
demostración, contrastar el texto original de un resultado o continuar una
línea de investigación.

En cada etapa, las acciones \emph{Estudie aquí}, \emph{Resuelva} y
\emph{Avance cuando pueda} forman la ruta interna. Toda instrucción
\emph{Lea}, \emph{Lea primero} o \emph{consulte} que mencione un libro o un
archivo externo es una lectura complementaria opcional: su ausencia no impide
seguir el desarrollo del volumen ni interpretar sus resultados.

\begin{center}
\fbox{\parbox{0.93\linewidth}{
\textbf{Convención para lecturas complementarias.} En los cuatro libros cuyo índice fue
fotografiado, ``p.'' significa la página \emph{impresa}. En los archivos de las
bibliotecas locales, ``PDF'' significa la página física que muestra el visor,
contada desde la primera hoja del archivo. No son numeraciones
intercambiables. Devaney, Hilborn y Guckenheimer--Holmes están escaneados en
pliegos dobles; sus rangos de visor se indican como aproximados.
}}
\end{center}

Si se desea ampliar el estudio, las fuentes locales originalmente inventariadas
se encuentran en:
\begin{sloppypar}
\begin{itemize}
  \item \texttt{C:\textbackslash Users\textbackslash moren\textbackslash
        Desktop\textbackslash Sistemas Dinamicos y Caos};
  \item \texttt{C:\textbackslash Users\textbackslash moren\textbackslash
        Desktop\textbackslash Calculo Fraccionario}.
\end{itemize}
\end{sloppypar}

\subsection{Regla de fuentes: el volumen para estudiar, las fuentes para contrastar}

Este volumen desarrolla las definiciones generales, las demostraciones
necesarias, los ejercicios y las técnicas numéricas. Los libros citados ofrecen
ampliaciones opcionales. Una afirmación sobre un atractor oculto
concreto se apoya primero en el artículo que publicó el modelo, los parámetros,
la condición inicial y la prueba de cuenca; después se contrasta con la
reproducción de este reporte. Los volúmenes recopilatorios sirven como mapa de
familias, no como sustituto del artículo original.

\begin{center}
\begin{tabularx}{0.96\linewidth}{@{}L{3.1cm}Y Y@{}}
\toprule
Pregunta & Fuente que se lee primero & Fuente que completa la lectura \\
\midrule
¿Qué significa oculto? &
Artículos fundacionales y de síntesis
\cite{LeonovKuznetsovVagaitsev2011,LeonovKuznetsov2013,DudkowskiEtAl2016} &
Libros recopilatorios
\cite{PhamEtAl2018HiddenAttractors,PhamVolosKapitaniak2017Systems,WangKuznetsovChen2021Multistability}. \\
¿Cómo se localizó un caso? &
Artículo original del método y del sistema
\cite{KuznetsovEtAl2017,Danca2017,Wu2023ArctanChua} &
Derivación algebraica y expediente numérico del caso en este reporte. \\
¿Qué prueba una gráfica? &
Artículos de diagnóstico, cuencas y fronteras
\cite{datseris2022,DazaEtAl2016BasinEntropy,DazaEtAl2015Wada} &
Textos de caos para aprender el observable y sus hipótesis. \\
¿Qué aporta un candidato bibliográfico? &
Artículo que fija ecuaciones, parámetros, orden e inicialización
\cite{Danca2017HiddenFO,munoz2018,wang2020} &
Libro compilado para comparar familias y realizaciones. \\
\bottomrule
\end{tabularx}
\end{center}

\subsection{Orden lógico: de la dinámica a la localización}

Los nombres de las materias no fijan por sí solos un orden de aprendizaje.
En particular, no conviene comenzar por exponentes de Lyapunov, atractores
extraños o dimensión fractal sin saber antes qué son un flujo, una órbita, un
conjunto invariante y una cuenca. La dependencia usada en este volumen es
\[
\begin{aligned}
&\boxed{\text{EDO y álgebra lineal}}
\longrightarrow
\boxed{\text{sistemas dinámicos}}
\longrightarrow
\boxed{\text{topología y geometría local}}
\\[0.55em]
&\hspace{2.2cm}\longrightarrow
\boxed{\text{caos y cuencas}}
\longrightarrow
\boxed{\text{Caputo y memoria}}
\longrightarrow
\boxed{\text{localización algebraico--numérica}}.
\end{aligned}
\]

La topología y la geometría aparecen dos veces. En una primera pasada dan el
lenguaje mínimo de proximidad, continuidad, invariancia, tangencia y
transversalidad. En una segunda pasada proporcionan grado, continuación,
variedades invariantes, homoclinicidad y espacios funcionales. Esa doble pasada
evita estudiar teoría abstracta sin saber todavía qué problema dinámico debe
resolver.

\begin{longtable}{@{}L{2.2cm}L{3.2cm}L{4.7cm}L{3.3cm}@{}}
\toprule
Bloque & Objeto central & Qué debe poder hacer al terminar & Para qué se usa después \\
\midrule
\endfirsthead
\toprule
Bloque & Objeto central & Qué debe poder hacer al terminar & Para qué se usa después \\
\midrule
\endhead
Sistemas dinámicos &
Flujo, órbita, equilibrio y ciclo &
Pasar de $\dot x=f(x)$ al campo, su flujo local, el jacobiano y la primera
variación &
Linealización, Poincaré, continuación y clasificación de destinos. \\
Geometría y topología mínimas &
Espacio de fases, conjuntos invariantes y cuencas &
Reconocer tangencia, transversalidad, separatrices y cambios de coordenadas &
Ocultedad, superficies críticas, PP y borde de cuenca. \\
Caos &
Recurrencia, expansión y complejidad &
Separar acotación, recurrencia, atracción, caos y ocultedad &
Evitar que una FFT o un exponente positivo se conviertan en una conclusión
excesiva. \\
Cálculo fraccionario &
Gamma, Mittag--Leffler, Caputo y Volterra &
Declarar orden, terminal inferior, rama compleja, datos iniciales e historia &
Transferencia $W_q$, Matignon, ABM--PECE y continuación causal. \\
Topología funcional &
Estado--historia y operadores integrales &
Distinguir cuenca funcional de cuenca de reinicio y retorno puntual de retorno
funcional &
FPP, continuidad con memoria y pruebas fraccionarias coherentes. \\
Métodos de localización &
Semillas DF, Machado, continuación, PP/FPP y borde &
Derivar cada semilla, declarar sus hipótesis y medir su residuo &
Banco unificado y campañas comparables. \\
Evidencia numérica &
Convergencia, clasificación y sondeos &
Comparar paso, horizonte, solucionador, memoria, destino y vecindades de
equilibrios &
Afirmaciones limitadas y reproducibles de atracción, caos y ocultedad. \\
\bottomrule
\end{longtable}

\subsection{Secuencia interna obligatoria y apoyo bibliográfico opcional}

\subsubsection*{Etapa 1. Campo vectorial, solución lineal y estabilidad local}

\textbf{Lea primero.}
\begin{enumerate}
  \item \emph{Teoría geométrica de ecuaciones diferenciales I}
        \cite{JaurezEtAl2021TGEDI}: cap.~1,
        p.~1; cap.~3, pp.~53--80; cap.~4, pp.~93--155; cap.~6,
        pp.~245--292.
  \item \emph{Teoría geométrica de ecuaciones diferenciales II}
        \cite{JaurezEtAl2021TGEDII}: cap.~8,
        pp.~69--98; en particular existencia y unicidad, p.~73; primera
        variación, p.~74; dependencia del dato inicial, p.~92; y
        Grobman--Hartman para dinámica discreta, p.~94.
  \item Si falta repaso de EDO: \emph{Métodos topológicos en el estudio de las
        ecuaciones diferenciales no lineales}
        \cite{Amster2021MetodosTopologicos}, apéndice A, pp.~319--327.
  \item Para fijar la intuición antes de las pruebas: archivo
        \texttt{Strogatz - Nonlinear Dynamics and Chaos.pdf}
        \cite{Strogatz2015Nonlinear}, PDF~30--50 para flujo y estabilidad y
        PDF~140--212 para sistemas lineales y plano de fases.
\end{enumerate}

\textbf{Estudie aquí.} Las definiciones de campo, flujo, órbita,
linealización, equilibrio e hiperbolicidad en
\cref{sec:geo-proposito}.

\textbf{Resuelva.} Para un sistema lineal $\dot x=Ax$, demuestre que
$\varphi_t(x_0)=e^{tA}x_0$ y que
$\varphi_{t+s}=\varphi_t\circ\varphi_s$. Después calcule el jacobiano y el
espectro del equilibrio nulo de uno de los casos enteros.

\textbf{Avance cuando pueda} explicar por qué la unicidad produce la ley de
flujo y por qué un valor propio con parte real nula impide aplicar directamente
Hartman--Grobman.

\subsubsection*{Etapa 2. Espacio de fases, topología y derivada de Lie}

\textbf{Lea.} En \emph{Topología y geometría diferencial con aplicaciones a
la física} \cite{NahmadAchar2022Topologia}: coordenadas adaptadas y curvas
integrales, §§11.3--11.5,
pp.~124--125; derivadas de Lie de funciones y campos, §§11.6--11.10,
pp.~125--127; derivada de Lie de formas, §§11.12--11.14, pp.~128--130;
ejercicios §§11.15--11.18, pp.~131--133. El índice fotografiado no muestra los
capítulos 1--10: no se asignan páginas que no fueron verificadas.

\textbf{Estudie aquí.} Espacio de fases, variedad, tangente, invariancia,
cuenca, frontera de cuenca y cilindro del PLL en
\cref{sec:geo-proposito}.

\textbf{Resuelva.} Pruebe que $\mathbb R\times S^1$ es el espacio natural del
PLL y escriba una distancia que identifique $\theta$ con $\theta+2\pi$.
Calcule también $L_fV=\nabla V\cdot f$ para $V(x)=\|x\|^2/2$.

\textbf{Avance cuando pueda} distinguir el dibujo de una trayectoria, la
órbita como conjunto y el flujo como familia de aplicaciones.

\subsubsection*{Etapa 3. Órbitas periódicas, Poincaré y Floquet}

\textbf{Lea en este orden.}
\begin{enumerate}
  \item \emph{Teoría geométrica de ecuaciones diferenciales II}
        \cite{JaurezEtAl2021TGEDII}, cap.~9:
        rectificación y Poincaré, pp.~113--115; Floquet y monodromía,
        pp.~121--141; Poincaré--Bendixson, p.~150; índice de Poincaré,
        p.~155; haz tangente y campos en variedades, p.~176.
  \item Archivo \texttt{Strogatz - Nonlinear Dynamics and Chaos.pdf}
        \cite{Strogatz2015Nonlinear}:
        PDF~215--260 para ciclos y Poincaré--Bendixson; PDF~261--325 para
        bifurcaciones, Hopf, retorno y Floquet.
\end{enumerate}

\textbf{Estudie aquí.} Sección transversal, mapa de retorno, monodromía,
multiplicadores y estabilidad orbital en \cref{sec:geo-proposito}.

\textbf{Resuelva.} Ejercicios 7.3.1--7.3.3 del archivo de Strogatz
(PDF~250--251) y 8.4.3--8.4.4 (PDF~313--314). Construya después una sección
transversal para el PLL o MAVPD y explique qué representa un multiplicador de
módulo mayor que uno.

\textbf{Avance cuando pueda} explicar por qué un punto fijo de Poincaré
representa una órbita periódica, no un equilibrio del campo original.

\subsubsection*{Etapa 4. Grado, disparo y continuación}

\textbf{Lea.} En \emph{Métodos topológicos en el estudio de las ecuaciones
diferenciales no lineales} \cite{Amster2021MetodosTopologicos}: condiciones
periódicas, p.~67; índice de curvas,
pp.~74--82; método de disparo, p.~84; Brouwer y Poincaré, p.~99; cap.~3,
grado de Brouwer, pp.~103--145; problemas, pp.~153--155.

\textbf{Estudie aquí.} \Cref{sec:metodos-topologicos-existencia}: residual de
disparo, punto fijo, grado de Brouwer, homotopía y continuación.

\textbf{Resuelva.} Calcule el grado de $F(x)=x$ y de $F(x)=-x$ en la bola
unidad. Después escriba el residual $P(z)-z$ de un mapa de retorno y enumere
las hipótesis que faltan antes de atribuirle un grado.

\textbf{Avance cuando pueda} distinguir «existe un cero» de «el cero atrae» y
de «el atractor es oculto».

\subsubsection*{Etapa 5. Caos: mecanismo, diagnóstico y cuencas}

\textbf{Paso 5A: mecanismo dinámico.} Lea primero el archivo de Strogatz
\cite{Strogatz2015Nonlinear}.
\begin{itemize}
  \item PDF~326--371: Lorenz y geometría del flujo;
  \item PDF~372--421: mapas, sensibilidad y exponentes de Lyapunov;
  \item PDF~422--445: dimensión fractal;
  \item PDF~446--476: atractores extraños.
\end{itemize}
Resuelva 9.2.2--9.2.3 y 9.3.8--9.3.9 (PDF~366--368),
10.3.7--10.3.8 y 10.5.1--10.5.4 (PDF~413--417), y 12.1.7 (PDF~472).

\textbf{Paso 5B: del dibujo al diagnóstico.} En el ejemplar físico
\emph{Exploring Chaos} \cite{Davies2004ExploringChaos}, lea puntos fijos y
estabilidad, p.~19; órbitas numéricas frente a exactas, p.~41; Fourier,
Lyapunov, caos y ergodicidad, pp.~44--60; bifurcaciones, pp.~65--106; y,
antes de usar \emph{edge tracking}, sistemas bidimensionales y fronteras de
cuenca, pp.~107--148, en especial §4.7, p.~134. Para flujos, complete con
plano de fase, Poincaré y ecuaciones variacionales, pp.~184--193.

Después use el ejemplar físico \emph{Chaos and Time-Series Analysis}
\cite{Sprott2003ChaosTimeSeries}: flujos no caóticos, pp.~46--71; teoría de
sistemas dinámicos, pp.~72--103; exponentes de Lyapunov, pp.~104--126;
atractores extraños, pp.~127--158; bifurcaciones, pp.~159--183; propiedades de
series, pp.~213--242; predicción y ruido, pp.~243--272; dimensión fractal,
pp.~302--328. El apéndice A, pp.~417--440, es un banco de comparación, no una
lista de atractores ocultos certificados.

\textbf{Paso 5C: cuencas y mecanismo topológico.} Lea
\textit{Alligood, Sauer \& Yorke - Chaos An Introduction to Dynamical
Systems.pdf} \cite{AlligoodSauerYorke1996Chaos}: cuencas, PDF~146--165;
fronteras y dimensión,
PDF~181--209; Lyapunov y herradura, PDF~211--247; atractores y medidas,
PDF~248--289; EDO y conjuntos límite, PDF~290--375; caos en flujos,
PDF~376--415; variedades, homoclinicidad y crisis, PDF~416--463. Resuelva los
\emph{Challenges} 4 (PDF~199--201), 7 (PDF~330--334), 8
(PDF~363--364) y 10 (PDF~444--451).

\textbf{Para demostraciones topológicas,} use
\texttt{Devaney - An Introduction to Chaotic Dynamical Systems.pdf}
\cite{Devaney1989Chaotic}:
dinámica simbólica, conjugación y caos, PDF aprox.~29--33; homoclinicidad,
PDF aprox.~70; atractores, hiperbolicidad y variedades invariantes,
PDF aprox.~110--125; Hopf, PDF aprox.~129. El archivo está escaneado a dos
páginas por hoja.
Si prefiere continuar en papel, \emph{A First Course in Chaotic Dynamical
Systems} es el acompañamiento físico anterior a este texto más formal. Como la
fotografía sólo muestra el lomo, no se le asignan páginas.

\textbf{Estudie aquí.} Homoclinicidad, herradura, dinámica simbólica, entropía,
fronteras fractales y límites de los diagnósticos en
\cref{sec:geo-proposito}.

\textbf{Avance cuando pueda} separar cinco afirmaciones: acotación,
recurrencia, atracción, caos y ocultedad. Ninguna implica automáticamente las
otras cuatro.

\subsubsection*{Etapa 6. Gamma, Mittag--Leffler, Laplace y Caputo}

\textbf{Paso 6A: definiciones y límites de validez.} En el ejemplar físico
\emph{The Fractional Calculus} \cite{OldhamSpanier1974Fractional}, lea Gamma,
p.~16; definiciones y equivalencias, pp.~47--60; propiedades generales,
pp.~69--94, en particular composición, p.~89, y dependencia del límite
inferior, p.~90. Lea Laplace, p.~136, después de fijar la definición del
operador. Esta obra da el lenguaje histórico; para el PVI causal de Caputo se
usará Diethelm como referencia principal.

\textbf{Paso 6B: cálculo guiado.} Lea en el archivo de ejercicios
fraccionarios \cite{Oliveira2019SolvedFractional}.
\begin{enumerate}
  \item \texttt{Oliveira - Solved Exercises in Fractional Calculus.pdf}:
        Gamma y funciones especiales, PDF~30--41; Mittag--Leffler,
        PDF~81--124; Laplace, PDF~130--137; definiciones GL, RL, Caputo y
        Hadamard, PDF~179--189.
  \item Ejercicios de Mittag--Leffler 13--16 y 30--32: enunciados
        PDF~142--145 y soluciones PDF~148--172.
  \item Ejercicios RL--Caputo 12, 14, 15, 20, 25, 32, 33 y 40: enunciados
        PDF~191--194 y soluciones PDF~197--224.
\end{enumerate}

\textbf{Paso 6C: interpretación y formulación causal.} Acompañe con
\texttt{Herrmann - Fractional Calculus An Introduction for Physicists.pdf}
\cite{Herrmann2014Fractional}: Gamma/Mittag--Leffler, PDF~28--35; GL y
Riemann--Caputo, PDF~36--53; integrales, derivadas y semigrupo,
PDF~64--89. Resuelva 3.3 (PDF~52; solución PDF~423), 5.1--5.2
(PDF~89; soluciones PDF~428--429) y 6.1 (PDF~97; solución PDF~432).

Termine con el ejemplar físico \emph{The Analysis of Fractional Differential
Equations} \cite{Diethelm2010AnalysisFDE}: operadores RL y GL, pp.~13--47;
definición y propiedades de Caputo, pp.~49--65; y Mittag--Leffler,
pp.~67--73. No continúe a los sistemas antes de poder explicar por qué RL,
GL y Caputo no son nombres intercambiables.

\textbf{Estudie aquí.} Funciones especiales, integral de
Riemann--Liouville, derivada de Caputo, transformada de Laplace y criterio de
Matignon en \cref{sec:fundamentos}.

\textbf{Avance cuando pueda} obtener
\[
\mathcal L\!\left\{t^{\beta-1}E_{\alpha,\beta}(\lambda t^\alpha)\right\}
=\frac{s^{\alpha-\beta}}{s^\alpha-\lambda}
\]
y explicar qué datos iniciales aparecen en la transformada de Caputo.

\subsubsection*{Etapa 7. Historia, Volterra y dinámica de Caputo}

\textbf{Paso 7A: memoria e inicialización.}
\begin{enumerate}
  \item En el archivo de Herrmann \cite{Herrmann2014Fractional}: no localidad
        y memoria, PDF~108--126;
        ejercicio 8.1, PDF~126; solución, PDF~436.
  \item En \texttt{Das - Kindergarten of Fractional Calculus.pdf}
        \cite{Das2020Kindergarten}:
        ramas complejas y cortes para $s^q$, PDF~41--47; integración RL,
        PDF~84--93; inicialización, kernel y convolución, PDF~109 y
        PDF~119--135; relación RL--Caputo, PDF~156--165; historia,
        inicialización y Laplace, PDF~247--286; FDE por Laplace,
        PDF~292--310; sistemas multivariables, PDF~333--375.
\end{enumerate}

\textbf{Paso 7B: PVI, regularidad y estabilidad.} En Diethelm
\cite{Diethelm2010AnalysisFDE}, lea existencia, pp.~85--92; unicidad,
pp.~93--108; perturbación de datos, pp.~109--115; regularidad,
pp.~116--126; ecuaciones lineales, pp.~133--156; y estabilidad,
pp.~157--166. El criterio angular de Matignon se estudia además en su artículo
original \cite{Matignon1996}; no debe atribuirse al libro.

\textbf{Paso 7C: operador integral y cálculo numérico.} En
\emph{Métodos topológicos} \cite{Amster2021MetodosTopologicos}, lea función de
Green, p.~229; Schauder, pp.~245--267; y Leray--Schauder, pp.~273--300.
Para ABM fraccionario consulte el algoritmo de una ecuación en Diethelm,
pp.~195--210, y el caso multi-término, pp.~211--228. Después contraste el
algoritmo con el artículo predictor--corrector original
\cite{DiethelmFordFreed2002}.

\textbf{Estudie aquí.} Ecuación integral de Volterra, perfil de memoria,
semidinámica funcional, cuenca funcional y retorno sobre historias en
\cref{sec:geometria-topologia-caputo}.

\textbf{Resuelva.} Construya dos historias $\phi_1$ y $\phi_2$ con el mismo
valor terminal y pasados distintos. Explique por qué el mismo punto actual no
determina necesariamente el mismo futuro. Derive además la forma de Volterra
del PVI de Caputo.

\textbf{Avance cuando pueda} explicar por qué reiniciar desde $x(t_*)$
cambia el problema y por qué una sección puntual de Poincaré deja de ser un
mapa de estado en general.

\subsubsection*{Etapa 8. Metodología de localización}

Estudie, en este orden, la forma de Lur'e y la transferencia fraccionaria
(\cref{sec:lure_transferencia}); la función descriptiva de Machado
(\cref{sec:machado-fdf}); integración y continuación causal
(\cref{sec:integracion_continuacion}); puntos perpetuos
(\cref{sec:perpetual-points}); y localización geométrico--topológica
(\cref{sec:localizacion-geometrico-topologica}).

Antes de la primera derivación de Lur'e, lea en Slotine--Li
\cite{SlotineLi1991} estabilidad avanzada y realimentación, PDF~119--175;
al comenzar la función descriptiva, lea PDF~176--209. Vuelva al libro sólo
después de derivar aquí la integral de Fourier, para que la fórmula tabulada no
reemplace el cálculo.

\textbf{Lea los artículos al llegar al método correspondiente, no todos al
principio.} Las tres carpetas de artículos proporcionadas se depuraron por
DOI, título y revista. Los resúmenes automáticos de SciSpace, presentaciones,
CSV y reportes internos no se emplean como literatura primaria.

\begin{longtable}{@{}L{3.0cm}L{6.0cm}L{4.4cm}@{}}
\toprule
Pregunta de tesis & Artículos principales & Qué se toma de ellos \\
\midrule
\endfirsthead
\toprule
Pregunta de tesis & Artículos principales & Qué se toma de ellos \\
\midrule
\endhead
Definición y origen &
\cite{LeonovKuznetsovVagaitsev2011,Leonov2013HiddenAttractors,LeonovKuznetsov2013} &
Definición operacional, distinción autoexcitado--oculto y primer mecanismo de
localización en Chua. \\
Función descriptiva y continuación &
\cite{KuznetsovEtAl2017,Machado2015FDF,BarbosaMachado2002BacklashImpact} &
Balance armónico, construcción de semilla y alcance auxiliar de la potencia
fraccionaria de Machado. \\
Casos fraccionarios &
\cite{Danca2017,Wu2023ArctanChua,Danca2017HiddenFO,munoz2018,wang2020} &
Ecuaciones, parámetros, orden, inicialización y evidencia publicada de cada
familia; no se transfieren conclusiones entre parámetros. \\
Puntos perpetuos &
\cite{Prasad2015Perpetual,DudkowskiPrasadKapitaniak2017,NazarimehrSaediJafari2017} &
Definición, uso como semilla y contraejemplos a su suficiencia. \\
Curvas y superficies de localización &
\cite{GilmoreEtAl2010Connecting,GuanXie2021Connecting,GodaraEtAl2018CriticalSurfaces} &
Objetos geométricos que generan semillas; no certifican por sí solos la
atracción ni la ocultedad. \\
Cuencas y fronteras &
\cite{datseris2022,DazaEtAl2016BasinEntropy,DazaEtAl2015Wada} &
Clasificación finita, incertidumbre y pruebas de frontera bajo dominio y
resolución declarados. \\
Síntesis del campo &
\cite{DudkowskiEtAl2016,GuanXie2025} &
Taxonomía y comparación de métodos; para un caso se vuelve siempre al artículo
original. \\
\bottomrule
\end{longtable}

\textbf{Uso de los libros compilados.} Para clasificar familias, consulte
primero \emph{Systems with Hidden Attractors}
\cite{PhamVolosKapitaniak2017Systems}: equilibrios estables, pp.~21--35;
infinitos equilibrios, pp.~37--50; sin equilibrios, pp.~51--63; y circuitos,
pp.~79--102. Use después \emph{Nonlinear Dynamical Systems with Self-Excited
and Hidden Attractors} \cite{PhamEtAl2018HiddenAttractors} como catálogo
amplio, y el capítulo fraccionario, pp.~199--238, de
\emph{Chaotic Systems with Multistability and Hidden Attractors}
\cite{WangKuznetsovChen2021Multistability}. De cada capítulo elegido extraiga
su referencia original y complete la ficha de ecuaciones, parámetros, orden,
condición inicial y prueba de cuenca antes de incorporarlo como candidato.

\textbf{Resuelva antes de pasar a los casos.}
\begin{enumerate}
  \item Derive $W_q(s)=c^{\mathsf T}(s^qI-A)^{-1}b$ y declare la rama de
        $s^q$.
  \item Obtenga una función descriptiva desde su integral de Fourier; no
        empiece por la fórmula final.
  \item Demuestre por qué $Df(p)f(p)=0$ define un punto perpetuo entero y por
        qué no puede interpretarse como una regla de cadena de Caputo.
  \item Dibuje, en un mismo esquema, semillas de función descriptiva,
        Machado, continuación, PP/FPP y borde de cuenca.
\end{enumerate}

\textbf{Puente entre la teoría estudiada y cada método.}
\begin{longtable}{@{}L{3.0cm}L{3.6cm}L{3.9cm}L{2.8cm}@{}}
\toprule
Método & Desarrollo algebraico indispensable & Desarrollo dinámico o geométrico & Prueba numérica mínima \\
\midrule
\endfirsthead
\toprule
Método & Desarrollo algebraico indispensable & Desarrollo dinámico o geométrico & Prueba numérica mínima \\
\midrule
\endhead
DF centrada, sesgada y Machado &
Realización de Lur'e, resolvente, integral de Fourier, rama de potencia y cierre
DC--fundamental &
Interpretar la solución como predictor de semilla, no como prueba de atractor &
Residuo del cierre, ramas, simetría e integración desde la semilla. \\
Continuación &
Familia $f_\lambda$, parametrización y condición inicial transportada &
Persistencia local y pérdida posible de una rama o de su estabilidad &
Malla en $\lambda$, paso refinado y misma historia causal. \\
PP y FPP &
$Df(p)f(p)=0$, $f(p)\ne0$, singularidad de $Df$ y expansión Picard--Caputo &
Covariancia afín, región suave y diferencia entre punto y evento con historia &
Residuo escalado, perturbaciones y destino con memoria completa. \\
Semillas geométricas y borde &
Superficies críticas, curvas conectantes y distancias escaladas &
Cuencas, separatrices, métrica cilíndrica y posible tercer destino &
Clasificador calibrado, tres brackets y dos solucionadores. \\
Caos y ocultedad &
Equilibrios, jacobianos, espectros y regiones de sondeo &
Atracción, recurrencia y exclusión de contacto con vecindades de todos los
equilibrios &
Convergencia de observables, controles negativos y radios predeclarados. \\
\bottomrule
\end{longtable}

\subsubsection*{Etapa 9. Casos, resultados y campaña dinámica}

No lea un caso como una sucesión de gráficas. Use tres pasadas, siempre en el
mismo orden:
\begin{enumerate}
  \item \textbf{Pasada algebraica:} escriba el campo, calcule equilibrios,
        jacobiano, polinomio característico, simetrías, realización, transferencia
        y ecuaciones de semilla que correspondan.
  \item \textbf{Pasada numérica:} declare solucionador, paso, horizonte,
        tolerancias, escala, transitorio y clasificador; para Caputo añada orden,
        terminal inferior, inicialización e historia completa.
  \item \textbf{Pasada de evidencia:} identifique destino, convergencia entre
        presupuestos, recurrencia, diagnóstico de caos, sondeos de cuenca y
        alcance exacto de la conclusión.
\end{enumerate}

Lea primero los enteros: Chua, Kalman--Fitts, MAVPD y PLL. Sólo después lea los
fraccionarios: Chua no suave, Chua arctangente, Lorenz generalizado y
Rabinovich--Fabrikant. Para cada sistema complete una ficha con siete renglones:
espacio de fases, equilibrios, estabilidad local, ruta de semilla, contrato de
integración, destino observado y alcance de la prueba de cuenca. Termine con los
protocolos y la campaña TG0--TG8.

\textbf{Compuerta antes de aceptar un resultado.} Debe poder reproducir los
residuos algebraicos; explicar por qué la semilla es sólo un punto de partida;
mostrar estabilidad del destino al refinar el cálculo; y distinguir una prueba
finita de vecindades de una afirmación global sobre la cuenca.

\textbf{Pregunta de control.} Si una trayectoria tiene exponente de Lyapunov
positivo, ¿qué falta para afirmar que su conjunto límite es un atractor oculto?
La respuesta debe mencionar atracción, robustez numérica y vecindades de todos
los equilibrios.

\subsubsection*{Etapa 10. Profundización; no es prerrequisito}

\begin{itemize}
  \item \emph{Topología y geometría diferencial con aplicaciones a la
        física}: conexiones, pp.~135--142; geodésicas, pp.~143--149;
        torsión y curvatura, pp.~151--165; métrica pseudo-riemanniana,
        pp.~165--186; variedades simplécticas, p.~194.
  \item \emph{Teoría geométrica de ecuaciones diferenciales II}: formas
        normales y bifurcaciones, pp.~185--244.
  \item \textit{Guckenheimer \& Holmes - Nonlinear Oscillations Dynamical
        Systems and Bifurcations.pdf} \cite{GuckenheimerHolmes1983}: flujo e
        invariancia, PDF aprox.~6--39;
        variedad central y formas normales, PDF aprox.~65--89; promediación y
        Melnikov, PDF aprox.~90--119; herradura y entropía,
        PDF aprox.~120--149; Lorenz, Shilnikov y conexiones globales,
        PDF aprox.~150--184.
  \item \texttt{Hilborn - Chaos and Nonlinear Dynamics.pdf}
        \cite{Hilborn2000Chaos}: Poincaré y
        bifurcaciones, PDF aprox.~38--60; homoclinicidad, herradura y
        Lyapunov, PDF aprox.~61--80; entropía, dimensiones y reconstrucción,
        PDF aprox.~162--220.
  \item En Diethelm \cite{Diethelm2010AnalysisFDE}: ecuaciones multi-término,
        pp.~167--186, sólo después de dominar el caso conmensurable.
  \item En Oldham--Spanier \cite{OldhamSpanier1974Fractional}:
        semiderivadas, pp.~120--132; técnicas numéricas y series,
        pp.~148--164; ecuación de Abel, p.~192; difusión, pp.~207--218.
\end{itemize}

\subsection{Qué función cumple cada archivo restante}

\begin{longtable}{@{}L{5.3cm}L{8.3cm}@{}}
\toprule
Archivo & Cuándo consultarlo \\
\midrule
\endfirsthead
\toprule
Archivo & Cuándo consultarlo \\
\midrule
\endhead
\texttt{Falconer - Fractal Geometry Mathematical Foundations.pdf}
\cite{Falconer2014Fractal} &
Después de estudiar fronteras de cuenca: dimensión de caja, PDF~59--75;
Hausdorff, PDF~76--97; cálculo de dimensiones, PDF~98--114; sistemas
dinámicos, PDF~238--264. \\
\texttt{Falconer - Fractal Geometry (Solution Manual).pdf} &
Sólo después de intentar el texto: soluciones 1--3, PDF~1--21; 4--8,
PDF~22--39; 9--13, PDF~40--71; 14--18, PDF~72--95. \\
\texttt{Barnsley - Fractals Everywhere.pdf} &
Después de Falconer si se requieren IFS: métricas y topología, PDF~12--48;
contracciones, PDF~49--129; dinámica, PDF~130--185; dimensión,
PDF~186--219. \\
\texttt{Barnsley - SuperFractals.pdf} &
Futuro: topología, PDF~19--99; transformaciones, PDF~100--200; semigrupos,
PDF~201--323; atractores IFS, PDF~324--395. \\
\texttt{Lorenz - The Essence of Chaos.pdf} &
Después de Strogatz: panorama, PDF~9--30; historia y atractor,
PDF~116--165; excursiones matemáticas, PDF~188--207. \\
\texttt{Mandelbrot - La Geometria Fractal de la Naturaleza (Espanol).pdf} &
Consulta visual: índice, PDF~7--9; dimensión, PDF~33; atractores fractales,
PDF~277; Browniano fraccionario desde PDF~367. Para pruebas use Falconer. \\
\bottomrule
\end{longtable}

\subsection{Ejemplares físicos y páginas verificadas}

Cuatro títulos de la fotografía no tienen copia en las dos carpetas locales.
Sus intervalos se verificaron en índices editoriales; por ello, en las etapas
5--7 se usa la página \emph{impresa}, no la página de un visor:

\begin{center}
\begin{tabularx}{0.97\linewidth}{@{}L{4.8cm}L{3.2cm}Y@{}}
\toprule
Ejemplar sin PDF local & Etapa & Lectura inicial \\
\midrule
\emph{Exploring Chaos} \cite{Davies2004ExploringChaos} &
5 & pp.~19, 41--60 y 107--148. \\
\emph{Chaos and Time-Series Analysis} \cite{Sprott2003ChaosTimeSeries} &
5 & pp.~46--183 y 213--272. \\
\emph{The Fractional Calculus} \cite{OldhamSpanier1974Fractional} &
6 & pp.~16, 47--60, 69--94 y 136. \\
\emph{The Analysis of Fractional Differential Equations}
\cite{Diethelm2010AnalysisFDE} &
6--7 & pp.~13--73, 85--126, 133--166 y 195--228. \\
\bottomrule
\end{tabularx}
\end{center}

Los ejemplares físicos de Strogatz y Devaney sí tienen copia PDF local y se
leen mediante los intervalos de visor indicados arriba. \emph{A First Course in
Chaotic Dynamical Systems} también está en papel y en Drive, pero se conserva
como acompañamiento; la secuencia principal ya queda cubierta por Strogatz,
Alligood y el texto avanzado de Devaney.

\subsection{Límites bibliográficos que no deben confundirse con resultados}

Los tres PDF locales de cálculo fraccionario no derivan el criterio de
Matignon ni el esquema ABM--PECE. Esos resultados se justifican con sus fuentes
específicas en el cuerpo matemático del reporte. El cap.~9 del archivo de Das,
dedicado a variantes modificadas de cálculo fraccionario, queda fuera de la
ruta Caputo. Del mismo modo, los textos de fractales no sustituyen la teoría de
flujos, cuencas, estabilidad o bifurcación.
